\documentclass[11pt]{article}

\usepackage[a4paper,margin=1in]{geometry}
\usepackage{amsmath}
\usepackage{amssymb}
\usepackage{enumitem}
\usepackage{parskip}
\usepackage[hidelinks]{hyperref}

\setlist[itemize]{topsep=2pt,itemsep=1pt,leftmargin=1.4em}
\setlist[enumerate]{topsep=2pt,itemsep=2pt,leftmargin=1.7em}

\title{\vspace{-2em}\textbf{\LARGE Network Position and Academic Performance}\\[0.4em]
\large Adding Network Centrality to Smirnov \& Thurner (2017)}
\author{}
\date{ECC3841 Network Economics --- Research Brief}

\begin{document}
\maketitle
\vspace{-2.5em}

\section*{Introduction}

A student's grades depend on effort and ability. They also depend on the student's
friends. A friendship network gives access to information, advice, study partners, and
motivation. So where a student sits in the network may matter for grades. We ask
whether it does.

\section*{Research question}

Does a student's position in a friendship network affect their grades, on top of who
their direct friends are?

\section*{Background}

Two papers frame the project.

\textbf{Smirnov and Thurner (2017)} tracked a changing ``who-likes-who'' network of
Russian students. They found that the link between friends' grades is mostly
\emph{selection}: students pick friends with similar GPA, rather than being pulled up or
down by them. Their method looks at one number per student --- the average GPA of that
student's direct friends. It ignores the rest of the network.

\textbf{Calvó-Armengol, Patacchini and Zenou (2009; IZA DP 3859)} supply the missing
part. In their model, a student's effort is proportional to their \emph{Katz--Bonacich
centrality}, a measure of network position. On US data, more central students do better
in school. But their data are one snapshot in time. So they cannot tell whether
position raises grades, or good grades just attract more friends.

We combine the two. We take the centrality prediction to Smirnov and Thurner's data,
which covers many time points, and test it while holding friend-group GPA and past GPA
fixed.

\section*{The model}

Each student $i$ picks study effort $x_i$. Their payoff is
\[
u_i = a\,x_i \;-\; \tfrac{1}{2}x_i^{2} \;+\; \phi \sum_{j} g_{ij}\,x_i x_j .
\]
The first term is the gain from effort. The second is its rising cost. The third is the
social part: $g_{ij}=1$ if $i$ and $j$ are friends, and $\phi>0$ means a friend
studying more raises $i$'s own reward from studying.

Setting the derivative to zero gives each student's best choice:
\[
x_i = a + \phi \sum_{j} g_{ij}\,x_j .
\]
Study a baseline amount, plus more when your friends study more. Solving for everyone at
once:
\[
\mathbf{x} = a\,(\mathbf{I} - \phi \mathbf{G})^{-1}\mathbf{1} = a\cdot \mathbf{b}(g,\phi).
\]
So effort --- and grades --- track $\mathbf{b}(g,\phi)$, the Katz--Bonacich centrality.
It expands as
\[
b_i = 1 + \phi\,(\text{direct friends}) + \phi^{2}\,(\text{friends of friends})
+ \phi^{3}\,(\dots) + \cdots
\]
a count of everyone you can reach by paths of any length, with longer paths weighted
less.

\section*{What \texorpdfstring{$\phi$}{phi} controls}

$\phi$ sets how far influence spreads through the network.
\begin{itemize}
  \item Small $\phi$: only direct friends count. Centrality is almost the same as
        \emph{popularity} (in-degree).
  \item Large $\phi$ (near its ceiling): distant links count almost as much as close
        ones. Centrality now measures \emph{reach} --- being a bridge between many
        groups.
\end{itemize}
So $\phi$ is a dial from ``popularity'' to ``reach.'' We compute centrality at several
values of $\phi$ and see how the result changes along the dial.

The ceiling on $\phi$ is $1/\lambda_{\max}$, where $\lambda_{\max}$ is the largest
eigenvalue of $\mathbf{G}$; above it the model has no solution. $\lambda_{\max}$ is
larger in bigger, denser networks, so the ceiling differs across our 38 networks. We
set $\phi = c \cdot (1/\lambda_{\max})$ with $c$ from $0.05$ to $0.95$. Then $c$ means
the same thing in every network, and results can be compared.

\section*{Empirical strategy}

For each network we compute centrality at each $c$, then run
\[
\begin{aligned}
\mathrm{GPA}_{i,\,t+1} = {}& \beta_0 + \beta_1\,\mathrm{Katz}_{i,t}
  + \beta_2\,\mathrm{GPA}_{i,t}
  + \beta_3\,\overline{\mathrm{GPA}}^{\,\text{friends}}_{i,t} \\[2pt]
  & {} + \beta_4\,\mathrm{InDeg}_{i,t} + \beta_5\,\mathrm{OutDeg}_{i,t}
  + \gamma_{\text{group}} + \varepsilon_{i,t}.
\end{aligned}
\]
Each control has a job. Average friend GPA is Smirnov and Thurner's variable, so
$\beta_1$ is net of friend-group composition. In- and out-degree are raw popularity, so
$\beta_1$ is net of that too. The real question is whether centrality adds anything
beyond counting friends. We run this at every $c$. One test, at $c = 0.85$, is fixed in
advance as the headline. We also try betweenness and eigenvector centrality.

\section*{Which way does the arrow run?}

Does position raise grades, or do grades attract friends? One snapshot cannot tell. Our
data help in three ways:
\begin{enumerate}
  \item \textbf{Time order.} Position is measured before the GPA it predicts.
  \item \textbf{Past GPA is controlled.} We compare students who currently have the same
        grade.
  \item \textbf{Panel data.} We can remove fixed traits (ability, background) with
        student fixed effects, and we can measure the reverse arrow directly: does a
        rise in GPA bring new friends?
\end{enumerate}
What we still cannot rule out: something that changes a student mid-term and lifts both
grades and friend count. ``Likes'' are also a rough stand-in for real friendship. So
our claim is that position \emph{predicts} later grades, not that it \emph{causes}
them.

\section*{Predictions}

\textbf{From the model:} $\beta_1 > 0$, and centrality beats raw degree. The claim
fails if $\beta_1$ is not significant once degree is controlled, or if it flips sign
along the $\phi$ dial.

\textbf{Reading the sign of $\beta_1$:}
\begin{itemize}
  \item $\beta_1 > 0$: central students gain from more access to information and help.
  \item $\beta_1 = 0$: the gains and the time costs cancel.
  \item $\beta_1 < 0$: wide reach costs study time.
\end{itemize}

\textbf{From replicating Smirnov and Thurner:}
\begin{itemize}
  \item A friend's past GPA does not predict a student's future GPA, once own past GPA
        is controlled.
  \item New friendships show a smaller GPA gap than ending ones --- a sign of selection.
  \item Selection is stronger at university than in high school, where students can
        change their circle more easily.
\end{itemize}

\section*{Data and network}

\textbf{Data:} Smirnov \& Thurner (2017) replication files, Harvard Dataverse
(\href{https://doi.org/10.7910/DVN/SZA9YW}{doi:10.7910/\allowbreak DVN/\allowbreak
SZA9YW}). About 6{,}200 Russian students, grades plus ``like'' ties every 3 months.

\textbf{Nodes:} students. 655 in high school; 1{,}200--1{,}549 per university year.

\textbf{Edges:} directed ``like'' ties. $i \to j$ means $i$ liked $j$ at least once in a
3-month window. 2 to 14 snapshots per group, 38 in total. Only 24--27\% of ties go both
ways.

\textbf{Attributes:} GPA; in-degree, out-degree, Katz--Bonacich (several $\phi$),
betweenness, eigenvector; high school vs.\ university and year.

\section*{Related studies}

\begin{enumerate}
  \item \textbf{Calvó-Armengol, Patacchini \& Zenou (2009), ``Peer Effects and Social
        Networks in Education''} (IZA DP 3859). The model above, tested on AddHealth
        friendship data. Limit: one snapshot, so cause and effect are not separated.
  \item \textbf{Giulietti, Vlassopoulos \& Zenou (2020), ``Peers, Gender, and
        Long-Term Depression''} (IZA DP 13680). Peer depression in school predicts adult
        depression, for women. Uses reference groups, not friendship ties. Zenou teaches
        this course.
  \item \textbf{Smirnov \& Thurner (2017),} \emph{PLOS ONE} 12(8): e0183473. The paper
        and data this project replicates and extends.
\end{enumerate}

\end{document}
